\documentclass[11pt]{article}
% =========================================================
%  學生講義 共用前言（以 XeLaTeX 編譯）
%  需要套件：xeCJK, tcolorbox（其餘多在 BasicTeX 內）
% =========================================================
\usepackage[a4paper,margin=2.3cm]{geometry}
\usepackage{amsmath,amssymb}
\usepackage{xcolor}
\usepackage{graphicx}

% ---- 中文字型（macOS 系統字型）----
\usepackage{xeCJK}
\setCJKmainfont{Songti TC}      % 內文：宋體
\setCJKsansfont{Heiti TC}       % 標題/強調：黑體
\xeCJKsetup{AutoFakeBold=2.2, AutoFakeSlant=0.2}

% ---- 版面 ----
\setlength{\parindent}{0pt}
\setlength{\parskip}{0.55em}
\linespread{1.15}
\renewcommand{\baselinestretch}{1.15}

% ---- 顏色 ----
\definecolor{cBlue}{HTML}{1f6feb}
\definecolor{cGray}{HTML}{6b7280}
\definecolor{cGrayBg}{HTML}{f3f4f6}
\definecolor{cPurp}{HTML}{7a3ff2}
\definecolor{cPurpBg}{HTML}{f5f1ff}

% ---- 標註框 ----
\usepackage[most]{tcolorbox}

% 就地補充新概念（灰底）：\begin{補充框}{標題}
\newtcolorbox{補充框}[1]{
  enhanced, breakable, arc=2pt, boxrule=0.6pt,
  colback=cGrayBg, colframe=cGray,
  left=9pt,right=9pt,top=7pt,bottom=7pt,
  before upper={\textbf{\sffamily #1}\par\vspace{3pt}},
}

% 延伸 / 開放問題（紫色左邊條）：\begin{延伸框}{標題}
\newtcolorbox{延伸框}[1]{
  enhanced, breakable, arc=2pt, boxrule=0pt,
  colback=cPurpBg, colframe=cPurpBg,
  borderline west={3pt}{0pt}{cPurp},
  left=10pt,right=9pt,top=7pt,bottom=7pt,
  before upper={\textbf{\sffamily 延伸閱讀｜#1}\par\vspace{3pt}},
}

% ---- 圖：置中、底下小字說明 ----
% \fig{檔名}{寬度}{說明}
\newcommand{\fig}[3]{%
  \begin{center}
  \includegraphics[width=#2\linewidth]{#1}\\[2pt]
  {\small\color{cGray} #3}
  \end{center}}

% ---- 章節標題用黑體 ----
\newcommand{\h}[1]{\sffamily #1}


\begin{document}

\begin{center}
{\sffamily\Huge\bfseries 選物II-3 牛頓運動定律的應用}\\[3pt]
{\sffamily\large 普通高中 選修物理II（力學二）・學生講義}
\end{center}
\vspace{0.6em}

本章不引入新的定律，而是把已學過的牛頓三大運動定律，應用到四類常見的力學問題：\textbf{靜力平衡}（物體為何維持不動）、\textbf{重心與傾倒}（物體在何種條件下會傾倒）、\textbf{摩擦力的應用}（物體會不會滑動）、\textbf{一維碰撞}（碰撞後的速度與能量去向）。這是力學的綜合應用單元。

橋樑、招牌、樓梯、起重機、煞車系統等結構與機械的設計，第一步都是確認其平衡狀態與摩擦條件。本章把抽象的 $\vec F=m\vec a$ 應用到具體可觀察的問題上。

\medskip
本章主線：

\begin{center}
靜力平衡（力矩） $\rightarrow$ 重心與穩定性 $\rightarrow$ 摩擦力的應用 $\rightarrow$ 一維碰撞
\end{center}

\medskip
碰撞在本章作為綜合應用處理：以\textbf{動量守恆}求碰後速度，再以\textbf{動能}判斷彈性或非彈性。動量、能量的完整理論分別見 [[選物II-1 動量與角動量]] 與 [[選物II-2 功與能量]]，本章引用其結論，不重新推導。

\section*{\sffamily 3-1\quad 靜力平衡：不平移，且不轉動}

\subsection*{\sffamily A.\ 從質點到剛體}

先前解力學問題時，多把物體視為一個\textbf{質點（particle）}：所有的力都作用在同一點，只要合力為零，物體即靜止。這對於判斷「會不會移動」是足夠的。

但實際物體有大小、有形狀。一塊招牌、一座橋、一把尺，力作用在\textbf{不同位置}。此時即使合力為零，物體\textbf{仍可能轉動}：一支鉛筆平放桌上，用兩根手指在兩端往\textbf{相反方向}各推一個相同大小的力，合力為零，但鉛筆會原地轉動。

因此對於\textbf{有大小的物體（剛體，rigid body）}，平衡需要\textbf{兩個}條件：合力為零（不平移），且\textbf{合力矩為零（不轉動）}。

\begin{補充框}{為什麼剛體平衡要「兩個」條件？}
合力為零保證物體\textbf{不平移}（重心不被整體推走）；但有大小的物體還可能\textbf{原地轉動}。要它「既不移、也不轉」，須再加一條「合力矩為零」。所以剛體平衡需要\textbf{兩個獨立條件}。

\textbf{為什麼質點不需要第二個條件？}質點是沒有大小的點，所有力都作用在\textbf{同一點}上，不存在「不同作用點」造成的轉動問題——力臂全為零，力矩自動為零。所以質點平衡只剩「合力＝0」一條。一旦物體有大小、力作用在不同位置，「轉不轉」就成為一個獨立的問題，必須另外用「合力矩＝0」處理。

兩個條件互相獨立：
$$\sum \vec F = 0 \;\Longrightarrow\; \text{不平移（質心加速度為零）};\qquad
\sum \tau = 0 \;\Longrightarrow\; \text{不轉動（角加速度為零）}.$$
可以「不移卻在轉」（如上述鉛筆的兩端反向推），也可以「在移卻不轉」（被一個通過重心的力推著直線加速）。要靜止不動，兩條都須滿足。
\end{補充框}

\subsection*{\sffamily B.\ 力矩：使物體轉動的量}

要描述「一個力使物體轉動的程度」，僅看力的大小不足，還須看\textbf{力作用在離轉軸多遠的位置}。推門時推門邊較省力、推靠近門軸處較費力，即由此而來。這個量稱為\textbf{力矩（torque）}。

\begin{補充框}{什麼是力矩與力臂？}
\textbf{力矩}＝一個力使物體繞某轉軸轉動的程度。同樣大的力，作用點離轉軸越遠、方向越接近垂直於連線，力矩越大。力 $F$ 對某一支點（轉軸）的力矩定義為
$$\boxed{\;\tau = F\,d\;},\qquad d = r\sin\theta,$$
其中 $r$ 是轉軸到力作用點的距離，$\theta$ 是力與這條連線的夾角，$d$ 是\textbf{轉軸到力的作用線的垂直距離}，稱為\textbf{力臂（lever arm）}。

\textbf{為什麼力臂是「垂直距離」？}把一個斜向的力拆成兩個分量：
\begin{itemize}\setlength{\itemsep}{1pt}
\item \textbf{沿著連線}的分量（指向或背離轉軸）：只會把桿件往軸的方向拉或壓，\textbf{不產生轉動}。沿門板方向推門，門不會轉。
\item \textbf{垂直於連線}的分量 $F\sin\theta$：這才是使物體轉動的部分。
\end{itemize}
有效的轉動分量是 $F\sin\theta$，乘上距離 $r$，得 $\tau = rF\sin\theta = F\,(r\sin\theta)=F\,d$。把 $\sin\theta$ 歸給距離，即得「力臂＝垂直距離 $d=r\sin\theta$」。兩種觀點（有效分力 $\times$ 距離，或全力 $\times$ 力臂）結果相同。力矩是\textbf{向量}，本章只處理平面問題，以\textbf{正負號}代表轉向即可（習慣取\textbf{逆時針為正}）。單位：牛頓·公尺（$\text{N·m}$）。

\textbf{扳手與開門的共同原理：}在 $\tau=Fd$ 中，$F$ 與 $d$ 可互換。力臂 $d$ 越大，產生同樣力矩所需的力 $F$ 越小。長扳手（大 $d$）可用較小的力鬆開螺帽；開門推遠離門軸處（大 $d$）較省力。這就是\textbf{槓桿}：以長力臂換取較小施力。
\end{補充框}

\fig{t3-torque}{0.95}{圖 3-1：力矩 $\tau=Fd$。左——開門：施力點離門軸越遠（力臂 $d$ 越大）越省力；沿門板方向（指向門軸）的力力臂為零，不產生轉動。右——扳手：力與桿成 $\theta$ 角時，力臂 $d=r\sin\theta$，故 $\tau=rF\sin\theta=Fd$。}

\textbf{注意：}
\begin{itemize}\setlength{\itemsep}{1pt}
\item 力臂\textbf{不是}「支點到施力點的距離」，而是「支點到\textbf{力的作用線}的垂直距離」。力未垂直於桿件時，須先乘上 $\sin\theta$，否則會高估力矩。
\item 沿著「指向支點或背離支點」方向的力（$\theta=0^\circ$ 或 $180^\circ$）力矩為零——它只造成拉伸或壓縮，不使物體繞該支點轉。力大不代表力矩大。
\item 「省力」不等於「省功」。長扳手讓施力變小，但手須移動更長的距離，做的功相同。槓桿省的是\textbf{力}，不是\textbf{功}。
\end{itemize}

\subsection*{\sffamily C.\ 剛體平衡的兩條件與取矩}

\begin{補充框}{剛體平衡條件}
剛體靜力平衡的條件為，同時滿足
$$\sum \vec F = 0 \quad(\text{各方向分量分別為零})\qquad\text{且}\qquad \sum \tau = 0.$$
第二式對\textbf{任意一點}取矩都成立：物體真的平衡時，對任一點取矩結果都是零。既然支點可任選，就挑能消去未知力的那一點——把支點選在某個未知力的作用點上，該力的力臂為零、自動從方程式中消去，未知數隨之減少。
\end{補充框}

\fig{t3-equilibrium}{0.62}{圖 3-2：剛體平衡的範例。水平桿左端以鉸鏈接在牆上（支點 O），右端由繩張力 $T$ 斜拉，桿自重 $W=mg$ 作用於中點，右端另掛招牌重 $W_2$，鉸鏈反力 $R$ 作用於 O。對 O 取矩時，$R$ 的力臂為零、自動消去，可直接由 $\sum\tau=0$ 求出 $T$。}

\textbf{注意：}力偶（couple）——兩個大小相等、方向相反、不在同一直線上的力——其合力為零，但有淨力矩，是純粹的轉動來源（轉動方向盤、轉螺絲起子皆為力偶）。「力偶合力為零所以沒有效果」並不正確。

\section*{\sffamily 3-2\quad 重心與平衡的穩定性}

\subsection*{\sffamily A.\ 重心}

一個有大小的物體，地球對它\textbf{每一小塊}都施加重力。把這些分散的重力加總，效果等同於有一個合力 $W=mg$ 作用在某一個特定點上，這個點稱為\textbf{重心（center of gravity）}。在均勻重力場中，重心與\textbf{質心（center of mass）}重合（質心的嚴謹定義見 [[選物II-1 動量與角動量]]）。對\textbf{對稱、均勻}的物體，重心位於幾何中心：尺的中點、圓盤的圓心、長方體的體心。

\subsection*{\sffamily B.\ 三種平衡}

把物體稍微偏離原位再放開，其反應取決於\textbf{偏離後重心是升高、降低、還是不變}。

\begin{補充框}{平衡的穩定性}
\begin{itemize}\setlength{\itemsep}{1pt}
\item \textbf{穩定平衡（stable）}：稍微偏離後\textbf{重心升高}，會自動回到原位（碗底的彈珠、不倒翁）。
\item \textbf{不穩定平衡（unstable）}：稍微偏離後\textbf{重心降低}，會越偏越遠、不回原位（倒立的鉛筆、山頂的球）。
\item \textbf{隨遇平衡（neutral）}：偏離後\textbf{重心高度不變}，停在任意位置都平衡（平地上的球、平躺的鉛筆）。
\end{itemize}
綜合而言：\textbf{重心越低、底面越大，越穩定}。
\end{補充框}

\subsection*{\sffamily C.\ 傾倒判準}

放在地面上的物體（櫃子、車輛、堆疊的積木）會不會傾倒，有一個直觀的判準。

\begin{補充框}{重心鉛垂線與支撐底面}
從\textbf{重心}畫一條鉛垂線（重力的作用線）：
\begin{itemize}\setlength{\itemsep}{1pt}
\item 鉛垂線\textbf{落在支撐底面（base of support）之內} → 重力對「即將翻倒的那條底邊」產生\textbf{回正}的力矩，物體被壓回，\textbf{不倒}。
\item 鉛垂線\textbf{落到支撐底面之外} → 重力的力矩變為\textbf{繼續翻倒}的方向，物體\textbf{傾倒}。
\item 鉛垂線\textbf{恰好通過底面邊緣} → 臨界狀態。
\end{itemize}

\textbf{用力矩理解這條判準。}物體會繞底面的某條邊緣（pivot）翻倒。當重心鉛垂線落在 pivot 內側，重力對 pivot 的力矩方向為「把物體壓回平放」（回正力矩）；落到 pivot 外側時，力臂變號，力矩變為「把物體往外扳」（翻倒力矩）；恰好通過 pivot 時力臂為零。所以「鉛垂線出界就倒」是\textbf{重力力矩變號的幾何判準}，背後仍是力矩平衡。

\textbf{什麼是支撐底面？}不是物體的底面積，而是\textbf{所有接觸點連起來圍成的最小凸區域}。四腳桌是四腳圍成的矩形；雙腳站立的人是兩腳掌與其間的區域；三腳架是三點圍成的三角形（三點必共平面，故三腳架在不平地面上也不晃）。\textbf{底面越大、重心越低，鉛垂線越不易跑出界，越穩定}。

\textbf{均勻物體的臨界角。}均勻物體（重心在幾何中心）立在地面，底寬 $w$、高 $h$，傾斜到重心鉛垂線恰好通過底邊時翻倒，臨界角
$$\tan\theta_c = \frac{w/2}{h/2} = \frac{w}{h}.$$
底越寬（$w$ 大）、越矮（$h$ 小），$\theta_c$ 越大，越難翻。車輛底盤做得寬而低，即為增大此臨界角；重心高的車輛（貨車、休旅車）或超載堆高的貨車則較易翻覆。
\end{補充框}

\fig{t3-stability}{0.92}{圖 3-3：傾倒判準。左——傾斜不大，重心鉛垂線落在支撐底面之內，重力產生回正力矩（穩定）；右——傾斜過大，重心鉛垂線落到支撐底面之外，重力產生翻倒力矩（傾倒）。}

\textbf{注意：}
\begin{itemize}\setlength{\itemsep}{1pt}
\item 「重心一定在物體\textbf{內部}」並不正確。甜甜圈、L 形板、高舉雙手跳過橫桿的跳高者，重心都可能落在物體\textbf{外面}（空心或凹處）。
\item 穩定與否取決於「\textbf{重心相對支撐底面}的位置」，不取決於物體的重量。空的高櫃比裝滿的矮櫃容易倒。
\item 臨界之前物體都還站立，但越接近臨界越不穩，一點擾動即傾倒。穩定性有程度之分。
\end{itemize}

\section*{\sffamily 3-3\quad 摩擦力的應用：滑或不滑}

\subsection*{\sffamily A.\ 靜摩擦、最大靜摩擦、動摩擦}

摩擦力的基本定義在 [[選物I-4 牛頓運動定律]] 已建立。本節聚焦於\textbf{判斷物體滑動或不滑動}，這是工程上常見的問題。

\begin{補充框}{摩擦力三段}
兩接觸面間，正向力為 $N$：
\begin{itemize}\setlength{\itemsep}{1pt}
\item \textbf{靜摩擦力 $f_s$}：物體尚未滑動時，摩擦力會\textbf{自動調整大小}以抵抗外力，範圍為 $0 \le f_s \le \mu_s N$。
\item \textbf{最大靜摩擦力}：靜摩擦能提供的上限
$$\boxed{\,f_{s,\max} = \mu_s N\,}.$$
外力一旦超過它，物體開始滑動。
\item \textbf{動摩擦力 $f_k$}：滑動中，大小固定
$$\boxed{\,f_k = \mu_k N\,},$$
方向恆與\textbf{相對滑動方向相反}。通常 $\mu_k < \mu_s$。
\end{itemize}

\textbf{判斷滑或不滑的步驟：}
\begin{enumerate}\setlength{\itemsep}{1pt}
\item 先\textbf{假設不滑}（靜摩擦），由平衡（或 $\sum F=ma$）求出維持此狀態\textbf{所需的靜摩擦力} $f_{\text{需}}$。
\item 算出\textbf{最大靜摩擦} $f_{s,\max}=\mu_s N$。
\item 比較：$f_{\text{需}} \le f_{s,\max}$ → \textbf{不滑}，實際摩擦力等於 $f_{\text{需}}$；$f_{\text{需}} > f_{s,\max}$ → \textbf{會滑}，改用動摩擦 $f_k=\mu_k N$ 重新計算。
\end{enumerate}

\textbf{注意：}「摩擦力＝$\mu N$」\textbf{只在物體正在滑動、或恰好將滑時才成立}。物體靜止且尚未達到極限時，靜摩擦力等於「維持平衡所需的值」，可遠小於 $\mu_s N$，外力為零時甚至為零。
\end{補充框}

\subsection*{\sffamily B.\ 斜面上的滑動判定}

把物體放在傾角 $\theta$ 的斜面上：重力沿斜面分量為 $mg\sin\theta$、垂直斜面分量為 $mg\cos\theta$（分解圖見 [[選物I-4 牛頓運動定律]]）。

\begin{補充框}{斜面臨界角}
物體在斜面上\textbf{恰好將滑}時，下滑分量等於最大靜摩擦：
$$ mg\sin\theta = \mu_s\,mg\cos\theta \;\Longrightarrow\; \boxed{\tan\theta_c = \mu_s}. $$
\begin{itemize}\setlength{\itemsep}{1pt}
\item $\theta < \theta_c$：不滑（靜止）。此時實際靜摩擦力等於下滑分量 $mg\sin\theta$，\textbf{不是} $\mu_s mg\cos\theta$。
\item $\theta > \theta_c$：下滑，加速度 $a = g(\sin\theta - \mu_k\cos\theta)$。
\end{itemize}
量出物體恰好開始滑動的角度，即可由 $\tan\theta_c=\mu_s$ 直接測出 $\mu_s$。
\end{補充框}

\textbf{靜摩擦與動摩擦在煞車上的差異。}輪胎鎖死打滑時用的是動摩擦 $\mu_k$。但 $\mu_s > \mu_k$，\textbf{將滑未滑}的滾動輪胎能提供較大的靜摩擦。防鎖死煞車系統（ABS）讓輪胎維持在「將滑未滑」邊緣，避免進入鎖死打滑，藉此維持較大的摩擦力、縮短煞車距離。另須注意，煞車距離 $d=v_0^2/(2\mu_k g)$ 正比於\textbf{速度平方}：速度加倍，煞車距離變為四倍。

\section*{\sffamily 3-4\quad 一維碰撞的綜合應用}

\subsection*{\sffamily A.\ 動量守恆與能量分析}

碰撞瞬間，兩物體間的作用力極大、時間極短。但這對力是\textbf{內力}（作用力與反作用力），對「兩物體系統」而言互相抵消；只要外力（重力、正向力等）在運動方向可忽略，系統總動量即守恆。

\begin{補充框}{碰撞的核心原理}
\textbf{任何碰撞}（不論彈性與否），只要外力可忽略，\textbf{系統總動量守恆}：
$$\boxed{\;m_1 v_1 + m_2 v_2 = m_1 v_1' + m_2 v_2'\;}$$
（一維，向右為正，撇號代表碰後。動量守恆的來歷見 [[選物II-1 動量與角動量]]。）

但\textbf{動能不一定守恆}。依碰後動能保留的多寡，把碰撞分成：
\begin{itemize}\setlength{\itemsep}{1pt}
\item \textbf{彈性碰撞（elastic）}：動能也守恆（理想情形，如剛性球、原子）。
\item \textbf{非彈性碰撞（inelastic）}：動能有損失（變形、發熱、發聲）。
\item \textbf{完全非彈性（perfectly inelastic）}：碰後\textbf{黏在一起}以共同速度前進，動能損失最多。
\end{itemize}

\textbf{動量與動能由不同機制保證。}動量守恆來自牛頓第三定律：碰撞時兩物互推的力大小相等、方向相反，對系統而言是內力、互相抵消，故不論碰撞多劇烈、物體是否變形發熱，總動量都守恆。動能（$\tfrac12 mv^2$）則是\textbf{巨觀整體運動}的能量；碰撞時物體可能變形、振動、發熱、發聲，這些過程把一部分巨觀動能轉成物體內部的能量。因此動能並未「消失」，而是\textbf{轉成其他形式的能量}（總能量永遠守恆）。
\end{補充框}

\fig{t3-collision}{0.86}{圖 3-4：兩種一維碰撞（B 原本靜止）。上——等質量彈性碰撞，動量與動能皆守恆，結果為 A 停止、B 以原速 $v$ 前進（交換速度）。下——完全非彈性碰撞，碰後 A、B 黏在一起以共同速度前進，動量守恆但動能有損失（轉為熱、聲、形變）。}

\subsection*{\sffamily B.\ 完全非彈性碰撞與彈性碰撞}

\textbf{完全非彈性碰撞}：碰後兩物以\textbf{共同速度} $v'$ 前進，未知數只有一個，僅用動量守恆即可求解。其動能損失最多，但\textbf{不會是全部}——動量守恆使黏在一起的整體仍須帶著總動量前進，這部分動能保留下來（僅在總動量為零，即等質量正面對撞時，動能才可能全部轉為內能）。

\textbf{彈性碰撞}：動量與動能皆守恆，兩個守恆方程式聯立可解出兩個碰後速度。以 $m_2$ 原本靜止為例：
\begin{itemize}\setlength{\itemsep}{1pt}
\item $m_1=m_2$：碰後 $v_1'=0$、$v_2'=v_1$，即\textbf{交換速度}（撞球母球停、子球前進）。
\item $m_1\gg m_2$：大撞小，$v_1'\approx v_1$、$v_2'\approx 2v_1$。
\item $m_1\ll m_2$：小撞大，$v_1'\approx -v_1$、$v_2'\approx 0$（原速反彈，如球撞牆）。
\end{itemize}

\begin{補充框}{彈性與非彈性的本質：動能去向}
把碰撞想成兩物之間夾著一根彈簧：\textbf{接近階段}彈簧被壓縮，動能轉為彈性位能；\textbf{分離階段}彈簧回彈，彈性位能轉回動能。
\begin{itemize}\setlength{\itemsep}{1pt}
\item 若這根彈簧是\textbf{理想}的（沒有能量漏成熱），釋放回來的動能與當初儲存的相同 → \textbf{動能守恆，彈性碰撞}。物體碰完不變形、不發熱，這是一種\textbf{理想化}。
\item 若彈簧不理想，甚至壓扁後不回彈（變形、發熱）→ 一部分能量留在內部 → \textbf{動能損失，非彈性碰撞}。
\end{itemize}
動能少掉的部分，等於變成熱、聲、形變的能量。實際上\textbf{沒有任何巨觀碰撞是完全彈性的}（總有少量發熱）；最接近彈性的是剛性鋼球，以及原子、分子、次原子粒子的碰撞（量子尺度常為嚴格彈性）。
\end{補充框}

\textbf{注意：}
\begin{itemize}\setlength{\itemsep}{1pt}
\item 「動量守恆代表動能也守恆」並不正確。動量在任何（外力可忽略的）碰撞都守恆，動能只在彈性碰撞守恆。
\item 「碰後黏在一起代表動量消失」並不正確。完全非彈性只是\textbf{動能}損失最多，\textbf{動量照樣守恆}（兩物一起帶著總動量前進）。
\item 「動能損失違反能量守恆」並不正確。損失的是\textbf{動能}這一項，它轉成了熱、聲、形變，\textbf{總能量守恆}。
\end{itemize}

\begin{延伸框}{恢復係數：量化彈性的程度}
一維碰撞的\textbf{恢復係數（coefficient of restitution）}定義為
$$ e = \frac{\text{碰後分離速率}}{\text{碰前接近速率}} = \frac{v_2' - v_1'}{v_1 - v_2}. $$
\begin{itemize}\setlength{\itemsep}{1pt}
\item $e=1$：完全彈性（動能守恆）。
\item $0<e<1$：一般非彈性（保留部分動能）。
\item $e=0$：完全非彈性（碰後同速，分離速率為零）。
\end{itemize}
$e$ 越接近 1，碰撞越接近彈性、保留的動能越多。籃球落地反彈 $e\approx0.8$（$e=\sqrt{h_{\text{反彈}}/h_{\text{落下}}}$），黏土球 $e\approx0$。碰撞段用動量守恆、與碰撞無關的運動段（如彈道擺的擺升）用力學能守恆，兩個守恆律可分段使用（力學能守恆見 [[選物II-2 功與能量]]）。本章僅要求對恢復係數的\textbf{定性理解}。
\end{延伸框}

\section*{\sffamily 3-5\quad 本章重點整理}

\begin{補充框}{一頁複習（公式與關鍵結論）}
\begin{itemize}\setlength{\itemsep}{2pt}
\item \textbf{力矩}：$\tau = Fd = rF\sin\theta$，$d$＝力臂（支點到力作用線的垂直距離）；逆時針取正。單位 N·m。
\item \textbf{剛體平衡兩條件}：$\sum\vec F=0$ 且 $\sum\tau=0$。取矩支點可任選，挑能消去未知力者（使其力臂為零）。
\item \textbf{重心與穩定性}：重心越低、底面越大越穩；\textbf{重心鉛垂線落出支撐底面即傾倒}。均勻物體臨界角 $\tan\theta_c=w/h$。
\item \textbf{摩擦三段}：靜摩擦 $0\le f_s\le\mu_s N$（自動調整）；最大靜摩擦 $f_{s,\max}=\mu_s N$；動摩擦 $f_k=\mu_k N$（通常 $\mu_k<\mu_s$）。判滑：比較「所需摩擦」與「$\mu_s N$」。
\item \textbf{斜面臨界角} $\tan\theta_c=\mu_s$；滑動後 $a=g(\sin\theta-\mu_k\cos\theta)$。煞車距離 $d=v_0^2/(2\mu_k g)\propto v_0^2$。
\item \textbf{碰撞}：總動量恆守恆 $m_1v_1+m_2v_2=m_1v_1'+m_2v_2'$；動能只在彈性碰撞守恆。完全非彈性碰後黏在一起（共速）、動能損失最多但動量仍守恆。恢復係數 $e=$（分離速率）/（接近速率），$e=1$ 為彈性、$e=0$ 為完全非彈性。
\end{itemize}
\end{補充框}

\end{document}
